\documentclass[10pt]{article}
%% korr-preamble.tex -- gemeinsame Präambel für die Korrekturlisten LNM 1184
%% Makros wie im Buch (LN-Book-Springer/preamble), damit die Formeln
%% genauso aussehen wie im Proof.
\usepackage[T1]{fontenc}
\usepackage[utf8]{inputenc}
\usepackage{newtxtext}
\usepackage[varvw]{newtxmath}
\usepackage[scaled=.9]{helvet}
\usepackage[varqu,scaled=.95]{zi4}
\usepackage[ngerman]{babel}
\usepackage{csquotes}
\usepackage[a4paper,landscape,margin=1.6cm,top=1.8cm,bottom=1.8cm]{geometry}
\usepackage{microtype}
\usepackage{array,longtable,booktabs}
\usepackage{ragged2e}
\usepackage[table]{xcolor}
\usepackage{xspace}
\usepackage{enumitem}
\setlist{nosep,leftmargin=*}
\usepackage[nointegrals]{wasysym}
\usepackage{newunicodechar}
\usepackage{nicefrac}
\usepackage{mathtools}
\usepackage{mleftright}
\mleftright
\usepackage[hidelinks]{hyperref}
\setlength{\parindent}{0pt}
\setlength{\parskip}{.5em}
\setlength{\emergencystretch}{3em}
\renewcommand{\arraystretch}{1.25}
\definecolor{codegray}{gray}{0.35}

%% ---- Makros aus dem Buch -------------------------------------------
\DeclarePairedDelimiterX{\norm}[1]{\lVert}{\rVert}{\ifblank{#1}{\:\cdot\:}{#1}}
\DeclarePairedDelimiterX{\abs}[1]{\lvert}{\rvert}{\ifblank{#1}{\:\cdot\:}{#1}}
\DeclarePairedDelimiterX{\dualp}[1]{\langle}{\rangle}{\ifblank{#1}{\,\cdot\,,\cdot\,}{#1}}
\newcommand{\interval}[1]{\left[ #1 \right]}
\newcommand{\ointerval}[1]{\left] #1 \right[}
\newcommand{\rointerval}[1]{\left[ #1 \right[}
\newcommand{\lointerval}[1]{\left] #1 \right]}
\newcommand{\tint}{\textstyle\int}
\AtBeginDocument{%
\let\ORGvarepsilon=\varepsilon
\let\epsilon=\epsilon
\let\epsilon=\ORGvarepsilon
\let\ORGvarrho=\varrho
\let\rho=\rho
\let\rho=\ORGvarrho
\let\ORGvartheta=\vartheta
\let\theta=\theta
\let\theta=\ORGvartheta
\let\ORGvarphi=\varphi
\let\phi=\phi
\let\phi=\ORGvarphi
\let\ORGvarleq=\leqslant
\let\leqslant=\leq
\let\leq=\ORGvarleq
\let\ORGvargeq=\geqslant
\let\geqslant=\geq
\let\geq=\ORGvargeq
\let\ORGvarpreccurlyeq=\preccurlyeq
\let\preccurlyeq=\preceq
\let\preceq=\ORGvarpreccurlyeq
\let\setminus\smallsetminus			% A - B fuer Mengen
\renewcommand{\P}{\mathfrak{P}}		% Achtung: \P hat andere Bedeutung in LaTeX
}
\newcommand{\N}{\mathbb{N}}		% Natuerliche Zahlen
\newcommand{\Z}{\mathbb{Z}}		% Ganze Zahlen
\newcommand{\Q}{\mathbb{Q}}		% Rationale Zahlen
\newcommand{\R}{\mathbb{R}}		% Reelle Zahlen
\newcommand{\C}{\mathbb{C}}		% Komplexe Zahlen
\newcommand{\K}{\mathbb{K}}		% Koerperzeichen
\newcommand{\T}{\mathbb{T}}		% Torus
\DeclareMathOperator{\Id}{Id}
\DeclareMathOperator{\sign}{sign}				% signum-Operator
\newcommand{\im}{\ensuremath{\mathrm{i}}}		% imaginäre Konstante i
\newcommand{\iu}{\ensuremath{\mathrm{i}}}
\newcommand{\eu}{\ensuremath{\mathrm{e}}}
\renewcommand*{\Re}{\mathop{}\!\mathrm{Re}}		% Realteil
\renewcommand*{\Im}{\mathop{}\!\mathrm{Im}}		% Imaginärteil
\newcommand*{\supp}{\mathop{}\!\mathrm{supp}\,}  % Träger
\newcommand*{\diff}[1]{\mathop{}\!\mathrm{d}{#1}}	% \diff{s} = ds, 
\newcommand*{\ds}{\mathop{}\!\mathrm{d}{s}}       	% \ds = ds, 
\newcommand*{\dg}{\mathop{}\!\mathrm{d}{g}}       	% \dg = dg, 
\newcommand*{\dt}{\mathop{}\!\mathrm{d}{t}}       	% \dt = dt, 
\newcommand*{\dx}{\mathop{}\!\mathrm{d}{x}}       	% \dx = dx, 
\newcommand*{\dy}{\mathop{}\!\mathrm{d}{y}}       	% \dy = dy, 
\newcommand*{\dr}{\mathop{}\!\mathrm{d}{r}}       	% \dr = dr, 
\newcommand*{\dm}{\mathop{}\!\mathrm{d}{m}}       	% \dm = dm, 
\newcommand*{\du}{\mathop{}\!\mathrm{d}{u}}       	% \du = du, 
\DeclareMathOperator{\esssup}{ess\,sup}
\newcommand{\LE}{\mathcal{L}(E)} 		% 
\newcommand{\BH}{\mathcal{B}(H)} 		% 
\newcommand{\BA}{\mathfrak{A}}
\newcommand{\ZE}{\mathcal{Z}(E)} 	% \ZE  Zentrum von E
\newcommand{\LH}{\mathcal{L}(H)} 
\let\originalL\L
\AtBeginDocument{%
  \DeclareRobustCommand{\L}[1]{\mathcal{L}(#1)}%
}
\newcommand{\Fix}[1]{\mathop{}\mathrm{Fix}\left(#1\right)}		% \Fix{T} Fixraum
\newcommand{\rank}{\mathop{}\!\mathrm{rank}\,}		% 
\newcommand{\Kern}[1]{\mathop{}\mathrm{ker}(#1)}
\newcommand{\Image}[1]{\mathop{}\mathrm{im}(#1)}
\newcommand{\image}{\mathop{}\!\mathrm{im}\,}
\renewcommand{\dim}{\mathop{}\!\mathrm{dim}\,}
\newcommand*{\eg}{e.g.,\xspace}
\newcommand*{\ie}{i.e.,\xspace}
\newcommand*{\resp}{resp.\xspace}
\newcommand*{\vs}{vs.\xspace}
\newcommand*{\cf}{cf.\xspace}
\newcommand*{\etc}{etc.\xspace}
\newcommand{\CA}{$\mathrm{C}^{*}$}	% C*-Algebra
\newcommand{\WA}{$\mathrm{W}^{*}$}	% W*-Algbera
\newcommand{\AW}{$\mathrm{AW}^{*}$}	% AW*-Algbera
\newcommand*{\TT}{\mathcal{T}}		% C_0-Semigroup
\newcommand*{\RR}{\mathcal{R}}		% Pseudo Resolvente
\newcommand*{\SG}{\mathcal{S}}	    % Semigroup
\newcommand*{\UG}{\mathcal{U}}	    % Implemented Semigroup
\newcommand*{\F}{\mathcal{F}}		% F-product
\newcommand{\mfrac}[2]{%
    {\raise.5ex\hbox{$\scriptstyle #1$}}%
    \kern.025em/\kern.025em%
    {\lower.25ex\hbox{$\scriptstyle #2$}}}
\newcommand*{\sfrac}[2]{\leavevmode\kern.1em%
       \raise.5ex\hbox{\small #1}%
       \kern.025em/\kern.025em%
       \lower.25ex\hbox{\small #2}}
\newcommand{\1}{\mathbf{1}}
\newcommand{\HO}{\mathcal{H}(O)} 		%  für Lotz
\newcommand{\WO}{\mathcal{W}(O)} 		%  für Lotz
\newcommand*{\ddp}{\mathop{}\!\mathrm{d}{p}}       	% \ddp = dp, 
\newcommand*{\dN}{\mathop{}\!\mathrm{d}{N}}       	% \dN = dN, 
\newcommand*{\ddt}{\mathop{}\!\mathrm\fraq{d}{dt}}   % \ddt = d/dt
\newcommand*{\lnm}{LNM1184: } % \LMN-BookCode
\newcommand{\closure}[1]{#1^{-}}
%% Hilfsmatrizen aus Teil I, Kap. 1
\newcommand{\ABC}{\begin{pmatrix} 0 & 1 \\ -1 & 0 \end{pmatrix}}
\newcommand{\BBC}{\begin{pmatrix} \cos(t) & \sin(t) \\ - \sin(t) & \cos(t) \end{pmatrix}}
\newcommand{\CBC}{\begin{pmatrix} 1 & 1 \\ -1 & -1 \end{pmatrix}}
\newcommand{\DBC}{\begin{pmatrix} 1+t & t \\ - t & 1-t \end{pmatrix}}
\newcommand{\EBC}{\begin{pmatrix} 0 & 1 \\ 1 &0 \end{pmatrix}}
\newcommand{\FBC}{\begin{pmatrix} \cosh(t) & \sinh(t) \\ \sinh(t) & \cosh(t) \end{pmatrix}}

%% ---- Darstellung der Listen ----------------------------------------
%% Quelltext, der sich nicht als Formel setzen lässt (Labels, \ref, Dateinamen)
\newcommand{\code}[1]{\leavevmode{\color{codegray}\ttfamily\small #1}}
\newcommand{\sichA}{\CIRCLE}      % ● hoch
\newcommand{\sichB}{\LEFTcircle}  % ◐ mittel
\newcommand{\sichC}{\Circle}      % ○ niedrig
%% Querverweise im Buch: Label-Namen statt Nummern
\newcommand{\lbl}[1]{\textsf{\footnotesize\color{codegray}\detokenize{#1}}}

\hypersetup{pdftitle={Mathematische Fehler in Beweisen: Part B (Teil II)}}
\pagestyle{plain}
\begin{document}
% Erzeugt aus math-error-part-b.md: Formeln gesetzt statt Quelltext.
% Grau/Schreibmaschine = Quelltext, der sich nicht als Formel setzen laesst (Makronamen, Labels).
{\Large\bfseries Mathematische Fehler in Beweisen: Part B (Teil II)\par}
\smallskip
\emph{One-parameter Semigroups of Positive Operators}, 2. Aufl. (Springer LNM 1184). Geprüft am 14.10.2026; Stand der Umsetzung: 20.10.2026.

\textbf{Umfang:} Beweise, Rechnungen und Beispiele in Teil II, Kap. 1–4, inhaltlich geprüft im LaTeX-Quelltext (Branch \code{fahnenkorrektur}): Vorzeichen und Konstanten, Ungleichungsrichtungen, Voraussetzungen und Quantoren, zitierte Hilfsaussagen, Lücken, Beispiele. Reine Satz- und Notationsfragen stehen in \code{korrekturen-mathe-part-b.md}; wo ein Befund dort schon steht und hier ergänzt wird, ist das vermerkt („Ergänzt korrekturen-mathe“).

\textbf{Wichtig:} Die Befunde sind von Claude nachgerechnet, aber nicht von einem zweiten Mathematiker geprüft. Bitte jeden Punkt selbst verifizieren, bevor er umgesetzt wird.

\textbf{Stand 20.10.2026:} Kap. 1 ist erledigt (nach Durchsicht durch Rainer Nagel umgesetzt, Branch \code{fahnenkorrektur}, als „Umgesetzt“ markiert). \textbf{Kap. 2, 3 und 4 sind noch zu bearbeiten.}

\textbf{Buchführung:} Maßgeblich für den Korrekturstand sind nur die Listen {korrekturen-text-*}, {korrekturen-mathe-*} und {math-error-*} (seit 20.10.2026). \code{Umgesetzte\_Korrekturen\_LNM1184.md}, \code{Finales-zum-Pruefen.md} und \code{Fahnenkorrektur\_Springer\_LNM1184.md} werden nicht mehr fortgeführt.

\textbf{Spalten:} S. $=$ gedruckte Buchseite (über das nächstgelegene Label ermittelt, $\ensuremath{\pm}1$ Seite). Z. $=$ Zeile in der .tex-Datei. Sicherheit: \sichA{} sicher falsch $\ensuremath{\cdot}$ \sichB{} sehr wahrscheinlich falsch $\ensuremath{\cdot}$ \sichC{} zweifelhaft, bitte prüfen.

Insgesamt \textbf{45 Punkte}: 28 \sichA{}, 9 \sichB{}, 8 \sichC{}.

\section*{Befunde}
\subsection*{Teil II, Kap. 1 – Basic Results on Spaces $C\ensuremath{_{0}}(X)$ (\code{part-b/chap-b1.tex})}
\begin{longtable}{@{}rrc>{\RaggedRight\arraybackslash}p{3.0cm}>{\RaggedRight\arraybackslash}p{11.4cm}>{\RaggedRight\arraybackslash}p{7.4cm}@{}}
\toprule
\textbf{S.} & \textbf{Z.} & \textbf{} & \textbf{Stelle} & \textbf{Befund} & \textbf{Vorschlag} \\
\midrule
\endfirsthead
\toprule
\textbf{S.} & \textbf{Z.} & \textbf{} & \textbf{Stelle} & \textbf{Befund} & \textbf{Vorschlag} \\
\midrule
\endhead
\bottomrule
\endlastfoot
139 & 139 & \sichC{} & Section 1.3, Adjungierte von Homomorphismen & \textbf{Umgesetzt.} {the adjoint of an algebra homomorphism induces a continuous map $\phi$ from $Y$ ... into $X$}: Für einen nicht-unitalen Algebrenhomomorphismus ist $T'\delta_{s}$ für manche $s$ gleich $0$ (z. B. $T=0$ oder $Tf = \1_{U}\cdot f(x_{0})$ mit offen-abgeschlossenem $U \subsetneq Y$), dann ist $\phi$ nicht auf ganz $Y$ definiert. Richtig nur für $T\1_{X} = \1_{Y}$ (so auch in (b) „Markov algebra homomorphism“ und stillschweigend in Kap. B-II, Z. 1645). & {the adjoint of a Markov algebra homomorphism} \\[2pt]
139 & 159 & \sichA{} & Section 1.3 (i), Darstellung von Verbandshomomorphismen auf $C(K)$ & \textbf{Umgesetzt.} {a function $h \in C(K)$}: Ohne Positivität von $h$ ist die „if“-Richtung falsch, z. B. $h \equiv -1$, $\phi = \mathrm{id}$ liefert $Tf = -f$, das kein Verbandshomomorphismus ist. Notwendig gilt $h = T\1 \geq 0$. & {a function $h \in C(K)_{+}$} \\
\end{longtable}
\subsection*{Teil II, Kap. 2 – Characterization of Positive Semigroups on $C\ensuremath{_{0}}(X)$ (\code{part-b/chap-b2.tex})}
\begin{longtable}{@{}rrc>{\RaggedRight\arraybackslash}p{3.0cm}>{\RaggedRight\arraybackslash}p{11.4cm}>{\RaggedRight\arraybackslash}p{7.4cm}@{}}
\toprule
\textbf{S.} & \textbf{Z.} & \textbf{} & \textbf{Stelle} & \textbf{Befund} & \textbf{Vorschlag} \\
\midrule
\endfirsthead
\toprule
\textbf{S.} & \textbf{Z.} & \textbf{} & \textbf{Stelle} & \textbf{Befund} & \textbf{Vorschlag} \\
\midrule
\endhead
\bottomrule
\endlastfoot
142 & 113 & \sichA{} & Beweis von Theorem 2.1.3, $(b) \ensuremath{\Rightarrow} (c)$ & Zu zeigen ist $(Af)(x) + \|A\|f(x) \geq 0$ nur für positive f (Aussage (c) ist $A + \|A\|\Id \geq 0$); „for all $f \in C_{0}(X)$“ ist falsch (für $-f$ statt f kehrt sich das Vorzeichen um), und der Beweis zeigt es in Z. 138–142 auch nur für $f \in C_{0}(X)_{+}$. & {for all $f \in C_{0}(X)_{+}$} \\[2pt]
146 & 401 & \sichB{} & Remark nach Lemma 2.1.9 (Gegenbeispiel zu Theorem 2.1.8 auf $C_{0}(0,1]$) & Mit $Af = f' + \frac{\alpha}{x}f$ und $f(0)=0$ liefert $(\lambda - A)f = g$ eindeutig $f(x) = -x^{-\alpha}\eu^{\lambda x}\int_{0}^{x}\eu^{-\lambda y}y^{\alpha}g(y)\dy$, also $R(\lambda,A) \leq 0$ (für $\ensuremath{\alpha} = 0$ ist das die Resolvente von $f \ensuremath{\mapsto} f'$ mit f(0)=0), nicht $\geq 0$. Mit umgekehrtem Vorzeichen der Ableitung gilt $R(\lambda,A)g(x) = x^{\alpha}\eu^{-\lambda x}\int_{0}^{x}\eu^{\lambda y}y^{-\alpha}g(y)\dy \geq 0$ für alle reellen $\ensuremath{\lambda}$ (ganz in $\ensuremath{\lambda}$, Bild in D(A)), und A ist kein Generator (Transport nach rechts mit Faktor $(x/(x-t))^{\alpha}$, unbeschränkt). & $Af(x) = -f'(x) + \frac{\alpha}{x} f(x)$; bitte mit [\lbl{arendt:1987}] abgleichen. \\[2pt]
147 & 491 & \sichA{} & Beweis von Proposition 2.1.10, $(a) \ensuremath{\Rightarrow} (b)$ & Wegen $\langle g,\phi_{x}\rangle = g(x)/u(x)$ ist $\langle A(-f),\phi_{x}\rangle = -u(x)^{-1}(Af)(x)$, nicht $-u(x)(Af)(x)$. Die Folgerung $(Af)(x) \geq 0$ bleibt richtig. & $-u(x)^{-1}(Af)(x) = \langle A(-f),\phi_{x} \rangle \leq 0$ \\[2pt]
149 & 656 & \sichC{} & Remark 2.1.15, letzter Satz & Als allgemeine Aussage falsch: $A = M_{m}$ mit $m(x) = -x$ auf C[0,1] erzeugt eine positive Halbgruppe, $s(A) = 0$, und $u = \1$ erfüllt $Au \leq s(A)u$ (Infimum in (2.1.8) angenommen), aber 0 ist kein Eigenwert ($mf = 0$ $\ensuremath{\Rightarrow} f = 0$). Unter der Voraussetzung „kompakte Resolvente“ vom Anfang der Remark ist s(A) ohnehin stets Eigenwert, der Satz wäre dann leer. & bitte prüfen (Satz streichen oder präzisieren, z. B. Voraussetzung/„eigenvalue of A'“ angeben) \\[2pt]
150 & 685 & \sichA{} & Beweis von Corollary 2.1.17 & Nach (2.1.4) ist $p_{u}(g) = \sup_{x} g^{+}(x)u(x)^{-1} \geq 0$ für alle g, also nie $p_{u}(-f) < 0$; „$f \ensuremath{\gg} 0$ iff $p_{u}(-f) < 0$“ ist falsch (es gilt nur $f \ensuremath{\geq} 0$ iff $p_{u}(-f) = 0$), und die Abschätzung $p_{u}(T(t)f) \leq \eu^{\lambda t}p_{u}(f)$ liefert nur die Positivität. Damit trägt der Beweis nicht. & Beweis ersetzen, z. B.: Wähle $u \gg 0$ in D(A) und $c \geq 0$ mit $Au \geq -cu$; dann $\eu^{ct}T(t)u - u = \int_{0}^{t}\eu^{cs}T(s)(A+c)u\ds \geq 0$, also $T(t)u \geq \eu^{-ct}u \gg 0$; für $f \gg 0$ gilt $f \geq \epsilon u$ und damit $T(t)f \geq \epsilon\eu^{-ct}u \gg 0$. \\[2pt]
150 & 723 & \sichA{} & Theorem 2.1.20, Gl. (2.1.10) und Beweis Z. 827 & (2.1.10) gilt nur für $\omega_{0}(A) \geq 0$: Für $A = -\Id$ und nicht konstantes $m \ensuremath{\gg} 0$ ist $\omega_{0}(m \cdot A) = -\min m > -\max m = \|m\|_{\infty}\omega_{0}(A)$. Im Beweis folgt $m \cdot Au \leq \|m\|_{\infty}\lambda u$ aus $Au \leq \lambda u$ nur für $\ensuremath{\lambda} \ensuremath{\geq} 0$. Für $\omega_{0}(A) < 0$ kehren sich die Schranken um (gleicher Beweis mit $\min m$ statt $\max m$). & (2.1.10) auf $\omega_{0}(A) \geq 0$ einschränken und ergänzen: für $\omega_{0}(A) < 0$ gilt $\|m\|_{\infty}\,\omega_{0}(A) \leq \omega_{0}(m \cdot A) \leq \|m^{-1}\|_{\infty}^{-1}\,\omega_{0}(A)$. \\[2pt]
152 & 900 & \sichA{} & Beweis von Example 2.1.22 (iii), auch Z. 916/917 & Die Abschätzung $\int_{-1}^{0}\eu^{\lambda x}\diff{\mu(x)} \leq \|\mu\|$ gilt nur für $\ensuremath{\lambda} \ensuremath{\geq} 0$; für $\ensuremath{\lambda} < 0$ kann $\lambda - \alpha - \int\eu^{\lambda x}\diff{\mu(x)}$ verschwinden, obwohl $\lambda > \alpha + \|\mu\|$. Gegenbeispiel: $\ensuremath{\alpha} = \ensuremath{-}10$, $\mu = \delta_{-1}$; dann ist $e_{\lambda}(x) = \eu^{\lambda x}$ für die Lösung $\ensuremath{\lambda} \ensuremath{\approx} \ensuremath{-}2,07$ von $\lambda + 10 = \eu^{-\lambda}$ Eigenfunktion, obwohl $\ensuremath{\lambda} > \ensuremath{-}9 = \ensuremath{\alpha} + \ensuremath{\|}\ensuremath{\mu}\ensuremath{\|}$. & {for $\lambda > \max\{0, \alpha + \|\mu\|\}$} (Z. 900, 916, 917) \\[2pt]
153 & 934 & \sichA{} & Gl. (2.1.14) und Z. 927 (nach Example 2.1.22) & Folgefehler zu (iii): $(\alpha + \|\mu\|,\infty) \subseteq \rho(A)$ und $s(A) \leq \alpha + \|\mu\|$ sind für $\ensuremath{\alpha} + \ensuremath{\|}\ensuremath{\mu}\ensuremath{\|} < 0$ falsch (Gegenbeispiel oben: $s(A) \ensuremath{\geq} \ensuremath{-}2,07 > \ensuremath{-}9$). & $(\max\{0,\alpha + \|\mu\|\}, \infty) \subseteq \rho(A)$ und $s(A) \leq \max\{0, \alpha + \|\mu\|\}$ \\[2pt]
162 & 1564 & \sichA{} & Example 2.2.13 & Die Einschränkung $\ensuremath{\alpha} \ensuremath{\geq} 0$ ist falsch: Der Beweis zeigt nur $\ensuremath{\beta} = 0$ ($\ensuremath{\alpha} \ensuremath{\in} \ensuremath{\mathbb{R}}$ beliebig), und für jedes $\ensuremath{\alpha} \ensuremath{\in} \ensuremath{\mathbb{R}}$ erzeugt A mit $L = \alpha\delta_{0}$ (nach Example 2.1.22 mit $\ensuremath{\mu} = 0$) die Halbgruppe $(T(t)f)(x) = f(x+t)$ bzw. $\eu^{\alpha(x+t)}f(0)$, die aus Verbandshomomorphismen besteht. & {if and only if\/ $L = \alpha\delta_{0}$ for some $\alpha \in \R$.} \\[2pt]
162 & 1628 & \sichA{} & Beweis von Example 2.2.13, Formel für T(t) & Mit $\eu^{t\alpha}f(0)$ (unabhängig von x) ist T(t)f an der Stelle $x = \ensuremath{-}t$ unstetig (linker Grenzwert f(0)), und die Randbedingung $f'(0) = \alpha f(0)$ wird nicht erhalten. Richtig ist die Fortsetzung durch die Lösung von $F' = \alpha F$. & $\eu^{\alpha(x+t)}f(0)   \text{if } x+t \geq 0$ \\[2pt]
173 & 2289 & \sichA{} & Example 2.3.15, Definition von $D_0$ (und Beweis Z. 2303) & Da F nur $\sup\|DF\| < \infty$ erfüllt, kann F linear wachsen; dann folgt aus $\mathrm{grad}\, f \to 0$ nicht $(\mathrm{grad}\, f \vert F) \in C_{0}$. Gegenbeispiel: $n = 1$, $F(x) = x$, $f = \sum_{n} n^{-1/2}b(x-n)$ (b glatte Beule): $f \in D_{0}$, aber $xf'(x)$ ist unbeschränkt, also $f \notin D(\delta)$. $D_0$ ist dann kein Kern, und die Behauptung $g \in C_{0}(\R^{n})$ in Z. 2303 ist unbegründet. & $D_{0} = \{ f \in C_{0}(\R^{n}) \cap C^{1}(\R^{n}) \colon (\mathrm{grad}\, f \vert F) \in C_{0}(\R^{n}) \}$ (invariant, da $(\mathrm{grad}(f\circ\phi_{t})\vert F) = (\mathrm{grad}\, f\vert F)\circ\phi_{t}$) oder $D_{0} = C_{c}^{1}(\R^{n})$; Invarianzbeweis Z. 2320–2372 entsprechend anpassen. \\[2pt]
175 & 2405 & \sichB{} & Definition 2.3.16 & Ergänzt korrekturen-mathe. Die Fallunterscheidung erfasst den Fall $c = a > \ensuremath{-}\ensuremath{\infty}$ (bzw. d $= b < \ensuremath{\infty}$) nicht, denn m(a) ist nicht definiert. Dann wäre z. B. $m \ensuremath{\equiv} 1$ auf (0,1) zulässig (keine Bedingung greift), obwohl $\ensuremath{\delta}_m$ kein Gruppengenerator ist, im Widerspruch zu Proposition 2.3.18 (dort: zulässig $\ensuremath{\Leftrightarrow} a' = \ensuremath{-}\ensuremath{\infty}$, $b' = \ensuremath{\infty}$) und zur Remark nach Theorem 2.3.23. Auch der Beweis von Theorem 2.3.17 (Z. 2658 „If $a < c$“) setzt $c = a$ als möglichen Fall voraus. & $c = a \text{ or } (a < c \text{ and } m(c) = 0)$ und analog $d = b \text{ or } (d < b \text{ and } m(d) = 0)$; bitte prüfen. \\[2pt]
175 & 2426 & \sichB{} & Satz nach Definition 2.3.16 & „every Lipschitz continuous function is admissible“ gilt nur, wenn an endlichen Randpunkten keine Bedingung verletzt wird: $m \ensuremath{\equiv} 1$ auf (0,1) ist Lipschitz-stetig, aber $\ensuremath{\delta}_m = d/dx$ auf $C_0(0,1)$ erzeugt keine Gruppe (Proposition 2.3.18, $a' = 0$). Richtig für (a,b) $= \ensuremath{\mathbb{R}}$ bzw. wenn m an endlichen Randpunkten gegen 0 geht. & {For example, every Lipschitz continuous function $m$ on $(a,b)$ with $\lim_{x \to a} m(x) = 0$ if $a > -\infty$ and $\lim_{x \to b} m(x) = 0$ if $b < \infty$ is admissible.} \\[2pt]
176 & 2630 & \sichA{} & Beweis von Theorem 2.3.17 (Kern-Aussage) & $T(t)f = f \circ \phi_{t}$ ist im Allgemeinen nicht differenzierbar: $\ensuremath{\varphi}_t$ ist in Nullstellen von m nicht notwendig differenzierbar. Gegenbeispiel: $m(x) = \vert x\vert$ auf $\ensuremath{\mathbb{R}}$ (Lipschitz, zulässig), $\phi_{t}(x) = x\eu^{t}$ für $x > 0$, $x\eu^{-t}$ für $x < 0$; für $f \in C_{c}^{1}(\R)$ mit $f(x) = x$ nahe 0 ist $f \in D_{0}(\delta_{m})$, aber $f \circ \phi_{t}$ ist in 0 nicht differenzierbar. Die Invarianz $T(t)D_{0}(\delta_{m}) \subseteq D_{0}(\delta_{m})$ ist also falsch, die Kern-Eigenschaft damit nicht bewiesen. & bitte prüfen (Kern-Eigenschaft anders begründen, z. B. durch direkte Approximation im Graphennorm) \\[2pt]
181 & 2931 & \sichB{} & Beweis von Theorem 2.3.24 & Für beliebige $C\ensuremath{^{1}}$-Diffeomorphismen $q_n$ folgt die Stetigkeit von $m = (q_{n}')^{-1}$ in Punkten von K nicht: $q_n'$ kann nahe einem endlichen Randpunkt $b_n$ immer wieder klein werden (nur $\int q_{n}' = \infty$ ist erzwungen), dann gilt $m(x) \not\to 0 = m(b_{n})$; ebenso an Häufungspunkten von Intervallen. Die $q_n$ müssen passend gewählt werden. & In Z. 2898 die Wahl präzisieren, z. B. $\lim_{x \to a_{n}} 1/q_{n}'(x) = \lim_{x \to b_{n}} 1/q_{n}'(x) = 0$ an endlichen Randpunkten und $\sup \vert 1/q_{n}'\vert \leq b_{n} - a_{n}$ (etwa $q_{n}(x) = \tan(\pi((x-a_{n})/(b_{n}-a_{n}) - 1/2))$ für beschränkte Intervalle). \\
\end{longtable}
\subsection*{Teil II, Kap. 3 – Spectral Theory of Positive Semigroups on $C\ensuremath{_{0}}(X)$ (\code{part-b/chap-b3.tex})}
\begin{longtable}{@{}rrc>{\RaggedRight\arraybackslash}p{3.0cm}>{\RaggedRight\arraybackslash}p{11.4cm}>{\RaggedRight\arraybackslash}p{7.4cm}@{}}
\toprule
\textbf{S.} & \textbf{Z.} & \textbf{} & \textbf{Stelle} & \textbf{Befund} & \textbf{Vorschlag} \\
\midrule
\endfirsthead
\toprule
\textbf{S.} & \textbf{Z.} & \textbf{} & \textbf{Stelle} & \textbf{Befund} & \textbf{Vorschlag} \\
\midrule
\endhead
\bottomrule
\endlastfoot
196 & 564 & \sichA{} & Beweis von Theorem 3.2.4 (i) & In der Kette $\tilde{T}\tilde{h}^{[n]} = \ldots = \lambda^{n} \cdot S_{\tilde{h}}^{n}\tilde{h} = \lambda^{n} \cdot \tilde{h}^{[n]}$ ist das dritte Glied falsch. Aus $\tilde{T}\vert\tilde{h}\vert = \vert\tilde{h}\vert$ folgt $\lambda^{n} S_{\tilde{h}}^{n}\vert\tilde{h}\vert$, und nach (3.2.7) ist dies $\lambda^{n}\tilde{h}^{[n]}$; $S_{\tilde{h}}^{n}\tilde{h}$ wäre dagegen $\tilde{h}^{[n+1]}$. & $= \lambda^{n} \cdot S_{\tilde{h}}^{n}\vert\tilde{h}\vert = \lambda^{n} \cdot \tilde{h}^{[n]}$ \\[2pt]
202 & 1036 & \sichC{} & Beweis von Korollar 3.2.12 & Zitiert wird Korollar 3.1.4. Es setzt $X$ kompakt voraus, hier ist $X$ aber beliebig lokalkompakt. Der Beweis von 3.1.4 braucht die Kompaktheit nur für $s(A) > -\infty$, und das liefert hier die Voraussetzung $\sigma(A) \neq \emptyset$. In Prop. 3.3.5 (v), Z. 1494, steht deshalb korrekt „see the proof of Corollary 3.1.4“. & {\ldots and the proof of Corollary~\lbl{cor:b3-1.4}} (analog Z. 1494) \\[2pt]
203 & 1096 & \sichB{} & Example 3.2.14 (i) & „by Corollary 3.2.10, there exists a dominant real eigenvalue with corresponding positive eigenfunction“: Kor. 3.2.10 liefert nur, dass $\sigma_{b}(A) = P\sigma_{b}(A)$ zyklisch ist. Daraus folgen weder die Dominanz (nach (3.2.19) könnte $\sigma_{b}(A)$ unbeschränkt sein) noch eine positive Eigenfunktion. Diese liefert Kor. 3.1.4, das aber $X$ kompakt voraussetzt; nötig ist hier $\sigma(A) \neq \emptyset$. & Begründung ergänzen, z. B. Dominanz über Kor.~\lbl{cor:b3-2.12} oder über die Reellwertigkeit der Eigenwerte (Eigenfunktionen liegen in $L\ensuremath{^{2}}$, dort ist A symmetrisch); positive Eigenfunktion „see the proof of Corollary~\lbl{cor:b3-1.4}“; bitte prüfen \\[2pt]
210 & 1551 & \sichB{} & Beweis von Theorem 3.3.6 (ii) & Aus (3.3.5), $A S_{h} = S_{h}(A + \im\alpha)$, folgt: $S_{h}$ bildet $\Kern{\im\alpha + A}$ auf $\Kern{A}$ ab. Das ist aber der Eigenraum zu $-\im\alpha$, nicht zu $\im\alpha$. Für die Behauptung über den Eigenraum zu $\im\alpha$ braucht man $S_{h}(\Kern{A}) = \Kern{\im\alpha - A}$, ebenfalls aus (3.3.5). Ohne den (nicht erwähnten) Schluss über $\bar{h}$ mit $A\bar{h} = -\im\alpha\bar{h}$ ist das Argument so nicht schlüssig. & {Assertion (iv) implies that $S_{h}$ maps $\Kern{A}$ onto $\Kern{\im\alpha - A}$ whenever \dots } \\[2pt]
211 & 1564 & \sichA{} & Beweis von Theorem 3.3.6 (i) & Ergänzt korrekturen-mathe. Bestätigt: $\Kern{A + \im(\beta - \alpha)} \neq \{0\}$ heißt $\Kern{\im(\alpha-\beta) - A} \neq \{0\}$, also $\im(\alpha - \beta) \in P\sigma(A)$. Für die Untergruppeneigenschaft genügt genau das, denn mit $\im\alpha, \im\beta$ liegt die Differenz in der Menge. $\im(\beta-\alpha)$ folgt erst zusätzlich aus der Symmetrie unter Konjugation, die im Beweis nicht gezeigt wird. & $\im(\alpha - \beta) \in P\sigma(A)$ \\[2pt]
211 & 1575 & \sichA{} & Beweis von Theorem 3.3.6 (iii) & „It follows from Proposition 3.3.5 (ii) that $h$ is strictly positive“: $h$ ist komplexwertig. Strikt positiv ist die Eigenfunktion $\vert h\vert$ (wegen $A\vert h\vert = 0$), und genau daraus folgt „h has no zeros“. & {that $\vert h\vert$ is strictly positive, \ie $h$ has no zeros in $X$} \\[2pt]
211 & 1595 & \sichA{} & Beweis von Theorem 3.3.6 (iii) & Ergänzt korrekturen-mathe. Bestätigt: $h\cdot p_{n}(\tilde{h})$ ist Linearkombination von $h^{[m]}$, $m \geq 1$, und $h^{[m]} \in \Kern{\im m\alpha - A}$ (Theorem 3.2.4 (ii) mit $A\vert h\vert = 0$). Mit $k = 0$ enthielte der Span $\Kern{A}$, auf dem $\phi$ nicht verschwindet ($\langle\vert h\vert,\phi\rangle > 0$). Der Schluss in Z. 1597–1598 ginge dann nicht mehr. & $k=1,2,3,\ldots$ \\[2pt]
214 & 1757 & \sichA{} & Example 3.3.10 (ii), Gl. (3.3.14) & In der Produktformel fehlt beim gemischten Term der Faktor 2: Für symmetrisches $(a_{ij})$ ist $\sum_{i,j} a_{ij}(u_{i}'w_{j}' + u_{j}'w_{i}') = 2\sum_{i}\sum_{j} a_{ij}u_{j}'w_{i}'$; ebenso Z. 1758 mit $(1/w)_{i}'$. Auf (3.3.15) hat das keinen Einfluss, weil sich die Terme erster Ordnung wegen $w_{i}' + w^{2}(1/w)_{i}' = 0$ ohnehin wegheben. & $\left(ub_{i} + 2\sum_{j} a_{ij}u_{j}'\right)w_{i}'$ bzw. $\left(ub_{i} + 2\sum_{j} a_{ij}u_{j}'\right)(1/w)_{i}'$ \\[2pt]
214 & 1769 & \sichC{} & Example 3.3.10 (ii) & „Since \dots{} $(a_{ij})$ is positive definite, we have grad $w = 0$“: $w = h/\vert h\vert$ ist komplexwertig. Aus $\sum a_{ij}w_{i}'w_{j}' = 0$ (ohne Konjugation) folgt mit positiver Definitheit allein nicht $w_{i}' = 0$; Beispiel: $\xi = (1,\im)$ bei $(a_{ij}) = \Id$. Der Schluss klappt erst wegen $\vert w\vert = 1$: Dann ist $w_{i}' = \im\theta_{i}w$ mit reellen $\theta_{i}$, also $\sum a_{ij}w_{i}'w_{j}' = -w^{2}\sum a_{ij}\theta_{i}\theta_{j}$. & Begründung ergänzen: „Since $\vert w\vert = 1$ we have $w_{i}' = \im\theta_{i}w$ with real $\theta_{i}$, hence \dots{}“; bitte prüfen \\[2pt]
222 & 2286 & \sichC{} & Example 3.4.7 (ii) & Die angegebene Menge $\{\hat{h}(r) + \im k\kappa(r)\} \cup \{h(z) \colon \kappa(\vert z\vert) = 0\}$ ist i. Allg. nicht abgeschlossen. Ist $\kappa(r_{0}) = 0$ und $\kappa(r) > 0$ für $r \to r_{0}$, so häufen sich die Punkte $\hat{h}(r) + \im k\kappa(r)$ gegen die ganze Gerade $\hat{h}(r_{0}) + \im\R$; ein Spektrum ist aber abgeschlossen. Vgl. den Abschluss in Theorem 3.4.11. & Abschluss der ersten Menge nehmen: $\overline{\{\hat{h}(r) + \im k\kappa(r) \colon \ldots\}}$; bitte prüfen \\[2pt]
223 & 2343 & \sichA{} & Beweis von Proposition 3.4.8 (iv) & Das angegebene $\phi_{k}$ ist i. Allg. kein Eigenvektor von $A'$: Der Integrand $\eu^{-\im 2\pi kt/\tau}h_{t}(x_{0})f(\phi(t,x_{0}))$ wächst pro Periode um den Faktor $\eu^{\tau\hat{h}(x_{0})}$. Damit gilt $\langle T(s)f,\phi_{k}\rangle = \eu^{(\hat{h}(x_0)+\im 2\pi k/\tau)s}\phi_{k}(f)$ nur für $\hat{h}(x_{0}) = 0$. Gegenbeispiel: konstantes $h = c \neq 0$ auf einer Kreisbahn, $f = e_{m}$, $m \neq k$. Erst mit dem zusätzlichen Faktor $\eu^{-\hat{h}(x_{0})t}$ wird der Integrand $\ensuremath{\tau}$-periodisch (vgl. das analoge $\phi_{nm}$ in (v)). & $\exp\left(-\left(\hat{h}(x_{0}) + \im\cdot\frac{2\pi k}{\tau}\right) t\right)\cdot h_{t}(x_{0})\cdot f(\phi(t,x_{0}))$ \\[2pt]
224 & 2365 & \sichA{} & Beweis von Proposition 3.4.8 (v), Gl. (3.4.13) & Nach der Rechnung davor ist $(\hat{h}+\im\beta - A')\phi_{nm} = \frac{1}{n}(\delta_{y} - c_{nm}\delta_{z})$ mit $y = \phi(m,x_{0}) \neq z = \phi(m+n,x_{0})$ (unendliche Bahn) und $\vert c_{nm}\vert \to 1$ für $m \to \infty$. Die Norm ist $(1 + \vert c_{nm}\vert)/n \to 2/n$, also stimmt $\leq 1/n$ nicht. Für den Schluss (Approximation mit $n \ensuremath{\to} \ensuremath{\infty}$) ist das unerheblich. & $\limsup_{m\to\infty} \|(\hat{h}+\im\beta -A')\phi_{nm}\| \leq 2/n$ \\
\end{longtable}
\subsection*{Teil II, Kap. 4 – Asymptotics of Positive Semigroups on $C\ensuremath{_{0}}(X)$ (\code{part-b/chap-b4.tex})}
\begin{longtable}{@{}rrc>{\RaggedRight\arraybackslash}p{3.0cm}>{\RaggedRight\arraybackslash}p{11.4cm}>{\RaggedRight\arraybackslash}p{7.4cm}@{}}
\toprule
\textbf{S.} & \textbf{Z.} & \textbf{} & \textbf{Stelle} & \textbf{Befund} & \textbf{Vorschlag} \\
\midrule
\endfirsthead
\toprule
\textbf{S.} & \textbf{Z.} & \textbf{} & \textbf{Stelle} & \textbf{Befund} & \textbf{Vorschlag} \\
\midrule
\endhead
\bottomrule
\endlastfoot
232 & 62 & \sichA{} & Theorem 4.1.1, Eq. (4.1.2) & $Af \leq \lambda f \text{ for some } 0 \leq f \in D(A)$: $f = 0$ erfüllt die Bedingung für jedes $\ensuremath{\lambda}$, das Infimum wäre also $\ensuremath{-}\ensuremath{\infty}$; auch für ein bloß positives $f \ensuremath{\neq} 0$ ist die Formel falsch. Die Charakterisierung gilt nur mit strikt positivem f, so wird sie auch im Beweis von Corollary 2.1.14 benutzt (chap-b2, Z. 635–640). & $\text{ for some strictly positive } f \in D(A)$; derselbe Fehler steht in chap-b2, Eq. (2.1.8), Z. 605. \\[2pt]
232 & 155 & \sichA{} & Example 4.1.2 (iii) & Für $\dot{r} = 1 - r$ ist X nur vorwärts invariant. $\ensuremath{\Phi}_t$ setzt sich stetig auf $\vert z\vert = 2$ fort, mit Bildern vom Radius $1+\eu^{-t}$, die im Inneren von X liegen. Deshalb verschwindet $f\circ\Phi_t$ für $f \ensuremath{\in} C_c(X)$ am Rand $\vert z\vert = 2$ im Allgemeinen nicht, d. h. T(t) bildet $C_0(X)$ nicht in $C_0(X)$ ab; außerdem erreicht r rückwärts in endlicher Zeit 0, es ist also kein Fluss auf $\ensuremath{\mathbb{C}}\ensuremath{\setminus}{0}$. & Bitte prüfen; z. B. $\dot{r} = r(1-r)(2-r)$ wählen, dann ist $\ensuremath{\Phi}_t$ ein Homöomorphismus von X. In (3), Z. 201, gilt dann allerdings nur noch punktweise $T(t)f \ensuremath{\to} 0$; mit majorisierter Konvergenz reicht das. \\[2pt]
235 & 353 & \sichA{} & Beweis von Theorem 4.1.4 & Part-I, Proposition 4.1.10 ($\ensuremath{\omega}_0 < 0 \ensuremath{\Leftrightarrow} \ensuremath{\|}T(t)\ensuremath{\|} \ensuremath{\to} 0 \ensuremath{\Leftrightarrow} \ensuremath{\|}T(t')\ensuremath{\|} < 1$) liefert aus $\int_0^{\infty} \|T(t)'\mu\| \dt < \infty$ nichts. Gebraucht wird der Satz von Datko, Part-I, Theorem 4.1.11, $(c) \ensuremath{\Rightarrow} (a)$ mit $p = 1$, angewandt auf die adjungierte Halbgruppe auf $C_0(X)^*$ (dort gilt die Integrierbarkeit für alle $\ensuremath{\mu} \ensuremath{\in} C_0(X)^* \ensuremath{\subseteq} C_0(X)'$). & {Part-I, Theorem~\lbl{thm:a4-1.11}} (Datko, $(c) \ensuremath{\Rightarrow} (a)$) statt {Part-I, Proposition~\lbl{prop:a4-1.10}} \\[2pt]
237 & 412 & \sichA{} & Theorem 4.2.1, Eq. (4.2.1) & $\|R_m(t)\| \leq C \cdot \exp((\epsilon +\Re(\lambda_m))t)$ ist zu schwach (dann wäre auch das $\ensuremath{\varepsilon}$ überflüssig). Für $m = 1$ ergibt sich in Corollary 4.2.2 nur $\|\eu^{-rt}T(t) - \nu\otimes h\| \leq C\eu^{\epsilon t}$, nicht die behauptete exponentielle Abnahme; dasselbe gilt für Theorem 4.2.5. Gemeint ist der nächste Eigenwert $\ensuremath{\lambda}_{m+1}$. & $\|R_m(t)\| \leq C \cdot \exp((\epsilon +\Re(\lambda_{m+1}))t)$ (mit $\Re(\lambda_{m+1}) \coloneq -\infty$, falls $\ensuremath{\sigma}(A)$ nur m Punkte hat) \\[2pt]
237 & 443 & \sichB{} & Beweis von Theorem 4.2.1 & Der Beweis liefert nur $\|T(t)(\Id -P)\| \leq \|\Id -P\|\cdot C_{0} \cdot \exp(\Re(\lambda_m) t)$ (Z. 447). Für die korrigierte Abschätzung des Rests braucht man $r(T(t_0)_{\vert \Kern{P}}) = \exp(\Re(\lambda_{m+1}) t_0)$ (Spektralabbildung für eventuell kompakte Halbgruppen). Daraus folgt, dass der Typ der Einschränkung höchstens $\Re (\ensuremath{\lambda}_{m+1})$ ist. & „the type of \dots{} is at most $\Re(\lambda_{m+1})$“ und in Z. 447 $C_{0} \cdot \exp((\Re(\lambda_{m+1})+\epsilon) t)$ \\[2pt]
237 & 455 & \sichA{} & Beweis von Theorem 4.2.1 & Ergänzt korrekturen-mathe. Mit $P_n E = \ker(\lambda_n-A)$ wäre $(A-\lambda_n)P_n = 0$, und die Summe in Eq. (4.2.1) würde sich auf $j = 0$ reduzieren. Das Bild des Residuums an einem Pol der Ordnung k(n) ist der verallgemeinerte Eigenraum; so wird es in Z. 458 ($(\lambda - A)^k = 0$) auch benutzt. & $P_n G = \ker(\lambda_n-A)^{k(n)}$ \\[2pt]
239 & 572 & \sichA{} & Proposition 4.2.3 & Ohne die Voraussetzung „D beschränkt“ ist $(c) \ensuremath{\Rightarrow} (a)$ falsch: Für $D = \{n\delta_x \colon n \in \N\}$ ist (c) erfüllt, weil x in höchstens einem $U_n$ liegt. D ist aber unbeschränkt und daher nicht relativ schwach kompakt. Der Satz von Grothendieck ist für beschränkte Mengen formuliert; in der Anwendung (Beweis von Prop. 4.2.4) ist D beschränkt. & „For a bounded subset $D \subseteq M(K)$ \dots{}“ \\[2pt]
247 & 1170 & \sichA{} & Beweis von Theorem 4.3.1 & $\|R(\lambda,B)\| \leq (\lambda-\omega_0)^{-1}$ gilt für die Wachstumsschranke $\ensuremath{\omega}_0$ von (S(t)) im Allgemeinen nicht, sondern nur $M(\lambda-\omega)^{-1}$ mit $M \ensuremath{\geq} 1$. Auch $\|R(\lambda,A_0)\| \leq (\lambda-\omega_0)^{-1}$ ist für $\ensuremath{\omega}_0 < 0$ falsch, weil $\ensuremath{\|}T_0(t)\ensuremath{\|} \ensuremath{\geq} 1$ für $t < 1$. Der Beweis funktioniert mit der Konstanten w aus der Standardannahme (B $\ensuremath{-} w$ erzeugt eine Kontraktionshalbgruppe, $w \ensuremath{\geq} 0$, Z. 1078), so wie in Z. 1408 und 1411. & In Z. 1139, 1148, 1170 und 1182 $\omega_0$ durch $w$ ersetzen; Z. 1139 „$\omega_0$ the growth bound of \dots{}“ entsprechend streichen. \\[2pt]
248 & 1193 & \sichB{} & Beweis von Theorem 4.3.1, Fall 1 & $h \colon [-t,0] \to F$, $h(r) \coloneq \delta_{r}T(-r)g$: Für $t > 1$ ist $\ensuremath{\delta}_r$ bei $r < \ensuremath{-}1$ nicht definiert. Gebraucht wird nur $r \ensuremath{\in}$ [s,0] (es gilt $s > \ensuremath{-}t$). & $h \colon [s,0] \to F$ und in Z. 1207 $r \in [s,0]$ \\[2pt]
249 & 1314 & \sichC{} & Remarks nach Corollary 4.3.2, (iii) & Der Schluss „u is not differentiable in t'“ setzt voraus, dass $\int_0^t S(t-s)\Phi(u_s)\ds$ in t' differenzierbar ist. Das folgt nicht, wenn $s \ensuremath{\mapsto} \ensuremath{\Phi}(u_s)$ nur stetig ist. Direkter: Eine Lösung erfüllt nach (ii) $u(0) =$ \ensuremath{\tilde{g}} $\ensuremath{\in} D(B)$, dann wäre $t \ensuremath{\mapsto}$ S(t)\ensuremath{\tilde{g}} überall differenzierbar – Widerspruch. & Bitte prüfen; Argument über Bedingung (ii) der Lösungsdefinition führen. \\[2pt]
251 & 1401 & \sichC{} & Text vor Proposition 4.3.5 & $w \in \R$: Die Abschätzung $\|\epsilon_\lambda \otimes R(\lambda,B)\Phi\| \leq (\lambda-w)^{-1} \cdot \|\Phi\|$ (Z. 1411) benutzt $\|\epsilon_\lambda\|_\infty = 1$, also $\ensuremath{\lambda} \ensuremath{\geq} 0$. Für $w < 0$ und $\ensuremath{\lambda} \ensuremath{\in}$ ($\ensuremath{\|}\ensuremath{\Phi}\ensuremath{\|} + w$, 0) ist $\|\epsilon_\lambda\|_\infty = \eu^{-\lambda} > 1$. & $w \in \R_{+}$ wie in Z. 1078 \\[2pt]
253 & 1537 & \sichC{} & Beweis des Lemmas (W. Arendt), Case 2 & „by the assumption $R(r,B+\Phi_{\mu}) \geq 0$“: Wegen $r \ensuremath{\leq} s(B+\ensuremath{\Phi}_\ensuremath{\mu})$ folgt aus $r \ensuremath{\in} \ensuremath{\varrho}(B+\ensuremath{\Phi}_\ensuremath{\mu})$ allein keine Positivität der Resolvente. Nötig ist der Grenzübergang $R(r,B+\Phi_{\mu+\epsilon}) \to R(r,B+\Phi_{\mu})$ in der Norm für $\ensuremath{\varepsilon} \ensuremath{\downarrow} 0$ (wegen $\ensuremath{\Phi}_{\ensuremath{\mu}+\ensuremath{\varepsilon}} \ensuremath{\to} \ensuremath{\Phi}_\ensuremath{\mu}$). & „\dots{} and, letting $\epsilon \downarrow 0$, $R(r,B+\Phi_{\mu}) = \lim R(r,B+\Phi_{\mu+\epsilon}) \geq 0$“ \\[2pt]
254 & 1569 & \sichA{} & Beweis von Corollary 4.3.8 & Theorem 4.3.7 (i) mit $\ensuremath{\lambda} = 0$ besagt: $s(B+\ensuremath{\Phi}_0) < 0 \ensuremath{\Rightarrow} s(A) < 0$, also $(b) \ensuremath{\Rightarrow} (a)$, nicht $(a) \ensuremath{\Rightarrow} (b)$. Für $(a) \ensuremath{\Rightarrow} (b)$: Ist $s(B+\ensuremath{\Phi}_0) = 0$, folgt $s(A) = 0$ aus Theorem 4.3.7 (ii); ist $s(B+\ensuremath{\Phi}_0) > 0$, folgt aus dem Lemma $s(A) = \ensuremath{\mu} \ensuremath{\geq} 0$. & Begründungen tauschen: „$(b) \ensuremath{\Rightarrow} (a)$ follows immediately from Theorem 4.3.7 (i)“; $(a) \ensuremath{\Rightarrow} (b)$ mit (ii) und dem Lemma beweisen. \\[2pt]
255 & 1582 & \sichA{} & Example 4.3.9 & $\alpha \in \C$: Corollary 4.3.8 (und Eq. (4.1.1)) setzt positive Halbgruppen voraus. Die Multiplikation mit $\ensuremath{\alpha}$ erzeugt auf $\ensuremath{\mathbb{C}}$ aber nur für reelles $\ensuremath{\alpha}$ eine positive Halbgruppe. Außerdem ist $s(B + \Phi_{0}) = \alpha + \|\Phi\|$ und die Bedingung $\alpha + \|\Phi\| < 0$ nur für reelles $\ensuremath{\alpha}$ sinnvoll. & $\alpha \in \R$ \\[2pt]
256 & 1669 & \sichB{} & Skalares Beispiel nach Example 4.3.10 (Claim) & $g(\lambda,\tau) \to p$ für $\ensuremath{\tau} \ensuremath{\to} \ensuremath{\infty}$ gilt nur für $\ensuremath{\lambda} > 0$; für $\ensuremath{\lambda} = 0 \ensuremath{\in} \ensuremath{\mathbb{R}}_+$ ist $g \ensuremath{\equiv} p + q$. Für den Zwischenwertschluss in $\ensuremath{\tau}$ bei festem $\ensuremath{\lambda}\ensuremath{{}'}$ braucht man außerdem $g(\ensuremath{\lambda}\ensuremath{{}'},\ensuremath{\tau}) \ensuremath{\to} p + q < \ensuremath{\lambda}\ensuremath{{}'}$ für $\ensuremath{\tau} \ensuremath{\downarrow} 0$, nicht $g(0,\ensuremath{\tau}) = p + q$. & „for every $\lambda > 0$“ und „$g(\lambda',\tau) \to p+q < 0$ as $\tau \downarrow 0$“ \\[2pt]
258 & 1773 & \sichA{} & Example 4.3.12 & $k \in C[a,1]$: L ist nur für $k \ensuremath{\geq} 0$ (auf [a,(r+1)/2]) positiv. Die Positivität von $\ensuremath{\Phi} = L(\ensuremath{\cdot}(\ensuremath{-}r))$ wird für Theorem 4.3.7 aber vorausgesetzt („It is easy to verify that $L$ is positive“). & $0 \leq k \in C[a,1]$ \\
\end{longtable}
\end{document}
